\begin{columns}[T]
\begin{column}{0.58\textwidth}
\footnotesize
\renewcommand{\arraystretch}{0.95}
\begin{tabular}{@{}lll@{}}
\toprule
\textbf{Model} & \textbf{Dropout*} & \textbf{Weight decay}\\
\midrule
Original transformer & 0.1 & 0\\
GPT2 & 0.1 & 0.1\\
T5 & 0.1 & 0\\
GPT3 & 0.1 & 0.1\\
T5 v1.1 & 0 & 0\\
PaLM & 0 & (variable)\\
OPT & 0.1 & 0.1\\
LLaMA & 0 & 0.1\\
Qwen 14B & 0.1 & 0.1\\
\bottomrule
\end{tabular}
\end{column}
\begin{column}{0.4\textwidth}
\vspace{2.5em}
\begin{itemize}
  \item Many older models used dropout during pretraining
  \item Newer models (except Qwen) rely only on weight decay
\end{itemize}
\end{column}
\end{columns}
\vfill
{\scriptsize * Most of the times papers just don't discuss dropout. On open models, this closely matches not doing dropout. This may not be true of closed models.\par}
